%! TeX program = lualatex
\documentclass[justified, nobib,a4paper]{tufte-handout}

\title{Linearized Euler equations}

\author[Wojciech Sadowski]{Wojciech Sadowski}

%\date{28 March 2010} % without \date command, current date is supplied

% \geometry{showframe} % display margins for debugging page layout

\usepackage{fontspec}
% \setmainfont[Renderer=Basic, Numbers=OldStyle, Scale = 1.0]{TeX Gyre Pagella}
% \setsansfont[Renderer=Basic, Scale=0.90]{TeX Gyre Heros}
% \setmonofont[Renderer=Basic]{TeX Gyre Cursor}

% Set up the spacing using fontspec features
\renewcommand\allcapsspacing[1]{{\addfontfeature{LetterSpace=15}#1}}
\renewcommand\smallcapsspacing[1]{{\addfontfeature{LetterSpace=10}#1}}

% setting ptions for included pictures
\usepackage{graphicx} % allow embedded images
\graphicspath{{graphics/}} % set of paths to search for images
\setkeys{Gin}{
  width=\linewidth,
  totalheight=\textheight,
  keepaspectratio
}

\hypersetup{colorlinks}

\usepackage[compatibility=false, font=footnotesize]{caption}
% * options are necessary to fix the bug that autoref was causing,
%   figure was constantly changed to section
\DeclareCaptionFont{blue}{\color{blue}}
\captionsetup{labelfont={blue}}

% color and color palettes
\usepackage{xcolor}

% drawings collor palette
\definecolor{p11}{HTML}{33cfff}
\definecolor{p12}{HTML}{ffa733}
\definecolor{p13}{HTML}{ff4133}


% drawing inside tex
\usepackage{tikz} 
\usetikzlibrary{decorations.markings}
\usetikzlibrary{calc}
\usetikzlibrary{shapes}

\usepackage{pgfplots}
\pgfplotsset{compat=newest}
\usepgfplotslibrary{fillbetween}
\usepackage{xcolor}

% math things
\usepackage{
  amsmath,  % extended mathematics
  amsthm,
  siunitx,  % non-stacked fractions and better unit spacing
  bm,       % bold math symbols
  physics2,  % a set of usefull symbols for differentiation 
  fixdif,
  derivative
}

\usephysicsmodule{ab, nabla.legacy}

% table things
\usepackage{
  booktabs, % book-quality tables
  multicol  % multiple column layout facilities
} 


% writing references
\usepackage{cleveref}  % must be loded after amsmath

% quoting text and friends
\usepackage{fancyvrb}         % extended verbatim environments
\fvset{fontsize=\normalsize}  % default font size for fancy-verbatim environments
\usepackage{
  dirtytalk, % provides \say command for easy quiting 
  nth       % provides \nth command for 1-st, 2-nd, etc
}        


% citing sources
\usepackage{natbib}
\setcitestyle{authoryear}
\bibliographystyle{plainnat}

% misc 
\usepackage{lipsum}   % filler text

\tikzset{
  set arrow inside/.code={\pgfqkeys{/tikz/arrow inside}{#1}},
  set arrow inside={end/.initial=>, opt/.initial=},
  /pgf/decoration/Mark/.style={
    mark/.expanded=at position #1 with
    {
      \noexpand\arrow[\pgfkeysvalueof{/tikz/arrow inside/opt}]{\pgfkeysvalueof{/tikz/arrow inside/end}}
    }
  },
  arrow inside/.style 2 args={
    set arrow inside={#1},
    postaction={
      decorate,decoration={
        markings,Mark/.list={#2}
      }
    }
  },
}

\input{../src/commands.tex}
\input{../src/symbols.tex}
\newtheorem{theorem}{Theorem}
\newtheorem{lemma}{Lemma}

\newcolumntype{L}[2]{>{\raggedleft#2}p{#1\textwidth}}


\newcommand{\Q}[1]{\section{Exercise:~{#1}}}
\newcommand{\Sol}{\subsection{Solution}}


\begin{document}

\maketitle% this prints the handout title, author, and date

\tableofcontents


\Q{Euler equations for isothermal flow (Laney)}
Show that for isothermal flow \(T=\const\), without the presence of body forces and energy sources, the Euler equations in one dimension
\[
	\begin{split}
		 & \pdv{\rho}{t} + \pdv{\rho u}{x} = 0, \\
		 & \pdv{\rho u}{t} + \pdv{\rho u^2}{x}
		= -\pdv{p}{x},                          \\
		 & \pdv{\rho E}{t} + \pdv{\rho E u}{x}
		=  -\pdv{u p}{x},
	\end{split}
\]
can be written as
\[
	\begin{split}
		 & \pdv{\rho}{t} + \pdv{\rho u}{x} = 0                     \\
		 & \pdv{\rho u}{t} + \pdv{}{x}\ab(\rho u^2 + \rho a^2) = 0
	\end{split}
\]
where the speed of sound is given as
\[
	a^2 = {\ab(\odv{p}{\rho})}_T = RT.
\]

\Sol{}
In general for ideal gas
\[
	p = (\gamma - 1) \rho e,
\]
and we know that \(e = c_v T\), with \(c_v=\const\), so
\[
	p = (\gamma - 1) \rho e = c_v(\gamma - 1)\rho T.
\]
Now we need to simplify the term with pressure in the momentum equation:
\[
	\pdv{p}{x_i} = \pdv{}{x_i}\ab[ c_v(\gamma - 1)\rho T] =
	c_v(\gamma - 1)T\pdv{\rho}{x_i} = {\ab(\odv{p}{\rho})}_T\pdv{\rho}{x_i}
	= \pdv{\rho a^2}{x_i}
\]
We can rearrange the momentum equation to get
\[
	\pdv{\rho u}{t} + \pdv{}{x}\ab(\rho u^2 + \rho a^2) = 0.
\]
The last thing: the energy equation:
\[
	\pdv{\rho E}{t} + \pdv{\rho E u}{x}  =  -\pdv{p u}{x},
\]
where the total energy is the sum of internal and kinetic energy, $E = e + u^2/2$.
Let's first simplify the tight side\sidenote{%
The derivative $\mdv{E}/{t}$ is the \emph{material} derivative, and is simply
equal to
\[
	\mdv{E}{t} = \pdv{E}{t} + \pdv{uE}{x}.
\]
We can also say that
\[
	\rho\ab(\pdv{E}{t} + \pdv{uE}{x}) =
	\pdv{E}{t} + \pdv{\rho uE}{x}
\]
using the continuity equation. Hence, we can write
\[
	\pdv{E}{t} + \pdv{\rho uE}{x} = \rho\mdv{E}{t}.
\]
}
\[
	\rho \mdv{E}{t} = \rho \mdv{}{t}\ab(c_v T + \frac{u^2}{2}) +
	\rho \ab( \underline{c_v \mdv{T}{t}} + \mdv{u^2 /2}{t} ) =
	\rho\mdv{u^2/2}{t}
\]
Of course, the underlined term has to be zero due to the assumption of
isothermal flow.
On the right,
\[
	-\pdv{p u_i}{x_i} = -\pdv{}{x_i}\ab[ c_v(\gamma - 1)\rho T u_i] =
	-a^2\pdv{\rho u_i}{x_i} = a^2\pdv{\rho}{t},
\]
so we end up with
\[
	\rho\mdv{u^2/2}{t} = a^2\pdv{\rho}{t}
\]
which just means that the change in density has to induce the change in kinetic
every of the gas. But we notice, that between, the equation of mass
conservation and newly derived momentum we have only two variables now. So we
essentially finished our exercise, and the total (or kinetic) energy is
essentially a \say{derived} quantity.

\Q{Linearized equations of gas dynamic}
\def\rs{\rho_\star}
\def\rz{\rho_0}
\def\us{u_\star}
\def\uz{u_0}
Provided, that the isothermal gas is moving with a uniform velocity \(\uz\) and a subject
to small disturbance \(\us\), has corresponding reference density \(\rz\) and
the density \(\rs\) associated with the flow disturbance, show that the linearized
equation of gas dynamics for this disturbance are
\begin{equation}
	\begin{split}
		 & \pdv{\rs}{t} + \pdv{\rz \us}{x} = 0                   \\
		 & \pdv{\us}{t} + \pdv{}{x}\ab(\rs \frac{a^2}{\rz}) = 0.
	\end{split}\label{eq:lin-eul}
\end{equation}
\subsection{Informal way}
We say \(u = \uz + \us\) and \(\rho = \rz + \rs\) and substitute here:
\[
	\begin{split}
		 & \pdv{\rho}{t} + \pdv{\rho u}{x} = 0                     \\
		 & \pdv{\rho u}{t} + \pdv{}{x}\ab(\rho u^2 + \rho a^2) = 0
	\end{split}
\]
so we have
\[
	\begin{split}
		 & \pdv{\rz + \rs}{t} + \pdv{\ab(\rz + \rs) \ab(\uz + \us)}{x} = 0                                                    \\
		 & \pdv{\ab(\rz + \rs) \ab(\uz + \us)}{t} + \pdv{}{x}\ab(\ab(\rz + \rs) {\ab(\uz + \us)}^2 + \ab(\rz + \rs) a^2) = 0,
	\end{split}
\]
and we do algebra. For the mass conservation
\[
	\pdv{\rz}{t} + \pdv{\rs}{t} + \pdv{\rz\uz}{x} + \pdv{\rz\us}{x} + \pdv{\rs\uz}{x} + \pdv{\rs\us}{x}= 0.
\]
We can directly remove all the terms involving only the derivatives of \(\rz\) and \(\uz\),
so we are left with:
\[
	\pdv{\rs}{t} + \pdv{\rz\us}{x} + \pdv{\rs\uz}{x} + \pdv{\rs\us}{x}= 0.
\]
Since both \(\rs\) and \(\us\) are disturbances, and we know that they are small, we can
simply neglect the last term on the left side, as the product \(\rs \us\) will be \say{doubly}
small.
So we are left with
\[
	\pdv{\rs}{t} + \pdv{\rz\us}{x} + \pdv{\rs\uz}{x} = 0.
\]
The momentum equation will be more tedious. The time derivative we treat in the same way
\[
	\pdv{\rho u}{t} = \pdv{\rs \uz}{t} + \pdv{\rz \us}{t},
\]
but the convection term is more of a mess,
\begin{fullwidth}
	\[
		\begin{split}
			\pdv{}{x}\ab(\rho u^2 + \rho a^2) =
			\pdv{}{x}\ab( \rz\uz^2 + 2\rz\uz\us + \rz\us^2  + \rs\uz^2 + 2\rs\uz\us + \rs\us^2 + \rz a^2 + \rs a^2),
		\end{split}
	\]
\end{fullwidth}
from which we remove all products of \(\star\) quantities and all derivatives of constant terms
\[
	\pdv{}{x}\ab(\rho u^2 + \rho a^2) =
	\pdv{}{x}\ab(  2\rz\uz\us + \rs\uz^2 + \rs\uz^2  + \rs a^2),
\]
So we are left with
\[
	\pdv{\rs \uz}{t} + \pdv{\rz \us}{t} + \pdv{}{x}\ab(  2\rz\uz\us + \rs\uz^2 + \rs a^2) = 0.
\]
Finally, we can make the following observation. If the fluid moves with the
uniform velocity \(\uz\), we can write our equations in a coordinate system
moving with velocity equal to \(\uz\). This will correspond to setting \(\uz=0\)
in our equations!\footnote{Euler equations should be invariant to Galilean
	coordinate change, prove it!} We can introduce this to be left with:
\[
	\begin{split}
		 & \pdv{\rs}{t} + \pdv{\rz \us}{x} = 0     \\
		 & \pdv{\rz\us}{t} + \pdv{\rs a^2}{x} = 0.
	\end{split}
\]
To get the form from the problem statement, simply divide the last equation by constant \(\rz\).

\subsection{Formal way}
Try it yourself:
\begin{enumerate}
	\item Write your equations as a vector valued function of \(\uz\) and \(\rz\),:
	      \[
		      \mathcal{F}(\uz, \rz) = \ab[
			      \begin{array}{c}
				      \pdv{\rz}/{t} + \pdv{\rz \uz}/{x} \\
				      \pdv{\rz \uz}/{t} + \pdv{\ab(\rz \uz^2 +\rz a^2 )}/{x}
			      \end{array}
		      ]
	      \]
	\item Linearize the equations by computing the generalized derivative
	      \[
		      \lim_{\epsilon \to 0} \frac{
			      \mathcal{F}(\uz + \epsilon \us, \rz + \epsilon \rs)
			      - \mathcal{F}(\uz, \rz)}{\epsilon}.
	      \]
	      This is the equation for a \say{tangent} to your original equation, but it is parameterized
	      by the reference values \(\uz\) and \(\rz\) that you haven't picked yet.
	\item Assume that the values are constant and \(\uz=0\), to get the final solution.
\end{enumerate}
For example for the continuity we will have:
\[
	\begin{split}
		 & \lim_{\epsilon \to 0}\frac{1}{\epsilon}\ab[
		\pdv{\rz + \epsilon \rs - \rz}{t} + \pdv{\rz \uz + \epsilon\rz\uz + \epsilon \rs \uz + \epsilon^2\rs\us - \rz\uz }{x}]=           \\
		 & \lim_{\epsilon \to 0}\frac{1}{\epsilon}\ab[
		\pdv{ \epsilon \rs}{t} + \pdv{ \epsilon\rz\uz + \epsilon \rs \uz + \epsilon^2\rs\us }{x}] =                                       \\
		 & \lim_{\epsilon \to 0}\pdv{ \rs}{t} + \pdv{ \rz\uz + \rs \uz + \epsilon\rs\us }{x} =\pdv{ \rs}{t} + \pdv{ \rz\uz + \rs \uz }{x}
	\end{split}
\]

\Q{Compute the flux Jacobian of equation for flow disturbance}
We are working with \cref{eq:lin-eul}, written now in matrix form
\[
	\pdv{}{t}
	\begin{bmatrix}
		\rs \\ \us
	\end{bmatrix}
	+ \pdv{}{x} \begin{bmatrix}
		\rz \us \\ a^2 \rs / \rz
	\end{bmatrix} = \bm{0}
\]
The variables are
\[
	\bm{u} = \begin{bmatrix}
		\rs \\ \us
	\end{bmatrix}
\]
and the flux function is:
\[
	\bm{f} = \begin{bmatrix}
		\rz \us \\ a^2 \rs / \rz
	\end{bmatrix}
\]
We know that we are looking for the form:
\[
	\pdv{}{t}
	\begin{bmatrix}
		\rs \\ \us
	\end{bmatrix}
	+ \pdv{\bm{f}}{\bm{u}}\pdv{}{x} \begin{bmatrix}
		\rs \\ \us
	\end{bmatrix} = \bm{0}
\]
Which is given as
\[
	\pdv{\bm{f}}{\bm{u}}\pdv{}{x} = \begin{bmatrix}
		0         & \rz \\
		a^2 / \rz & 0
	\end{bmatrix}
\]

\Q{Diagonalize the Jacobian (Is the system hyperbolic?)}
The eigenvalue problem is given as
\[
	\bm{A}\bm{k} = \lambda \bm{k}
\]
\[
	\bm{A}\bm{k} - \lambda \bm{k} = \bm{0}
\]
\[
	\ab(\bm{A} - \lambda \bm{I})\bm{k} = \bm{0}
\]
Equation has a non-zero solution if and only if the determinant of the matrix \(\bm{A} - \lambda \bm{I}\) is zero.
\[
	\det \ab(\bm{A} - \lambda \bm{I}) =\det \begin{bmatrix}
		-\lambda  & \rz      \\
		a^2 / \rz & -\lambda
	\end{bmatrix} = \lambda^2 - a^2 = \ab(\lambda - a)\ab(\lambda + a) = 0
\]
The diagonal matrix is equal to
\[
	\bm{\Lambda} = \begin{bmatrix}
		-a & 0 \\
		0  & a
	\end{bmatrix}
\]
\Q{Compute eigenvectors}
\[
	\bm{A}\bm{k}^i = \lambda^i \bm{k}^i
\]
or in full
\[
	\begin{bmatrix}
		-\lambda^i & \rz        \\
		a^2 / \rz  & -\lambda^i
	\end{bmatrix}
	\begin{bmatrix}
		k^i_1 \\
		k^i_2
	\end{bmatrix} = \bm{0}
\]
so if we write the actual equations for \(\lambda = -a\)
\[
	\begin{split}
		a k_1 + \rz k_2     & = 0  \stackrel{\times a/\rz}{\rightarrow}  a^2/\rz k_1 + a k_2 = 0 \\
		a^2/\rz k_1 + a k_2 & = 0
	\end{split}
\]
The equations are identical, so we will look for parametric solution.
It is clear that the solution is of the form
\[
	k^1 = \alpha\begin{bmatrix}
		1 \\
		-a/\rz
	\end{bmatrix}
\]
where \(\alpha\) is the parameter.
The second vector similarly is:
\[
	k^2 = \alpha\begin{bmatrix}
		1 \\
		a/\rz
	\end{bmatrix}
\]
We can pick \(\alpha = \rz\), so our vectors are:
\[
	k^1 = \begin{bmatrix}
		\rz \\
		-a
	\end{bmatrix},\quad
	k^2 = \begin{bmatrix}
		\rz \\
		a
	\end{bmatrix}.
\]
The system is hyperbolic, because it has two real eigenvalues, two different eigenvectors.

\Q{Derive the characteristic form}
We are looking for the equation of the form:
\[
	\pdv{\bm{w}}{t} + \bm{\Lambda}\pdv{\bm{w}}{x} = 0
\]
where, $\bm{w}$ is the solution is terms of characteristic variables. Transform Euler equations into the equations above.
\Sol
We know that:
\[
	\bm{A} = \bm{K}\bm{\Lambda}\bm{K}^{-1},
\]
where, the matrix of eigenvectors is
\[
	\bm{K} = \begin{bmatrix}
		k^1_1 & k^2_1 \\
		k^1_2 & k^2_2
	\end{bmatrix} = \begin{bmatrix}
		\rz & \rz \\
		-a  & a
	\end{bmatrix}.
\]
If we put that inside our equation we get:
\[
	\pdv{\bm{u}}{t} + \bm{K}\bm{\Lambda}\bm{K}^{-1}\pdv{\bm{u}}{x} = 0.
\]
We can propose a variable change so $\bm{u} = \bm{K}\bm{u}$, and put that into the
equation (remember that $\bm{K}$ has only constant coefficients)
\[
	\bm{K}\pdv{\bm{w}}{t} + \bm{K}\bm{\Lambda}\underbrace{\bm{K}^{-1}\bm{K}}_{=\bm{I}}\pdv{\bm{w}}{x} = 0.
\]
We can now extract the matrix $\bm{K}$
\[
	\bm{K}\ab(\pdv{\bm{w}}{t} + \bm{\Lambda}\pdv{\bm{w}}{x}) = 0
\]
and multiply by the inversion
\[
	\bm{K}^{-1}\bm{K}\ab(\pdv{\bm{w}}{t} + \bm{\Lambda}\pdv{\bm{w}}{x}) = 0
\]
which gets us where we need to be.

\Q{Derive the initial condition in terms of characteristic variables}
The initial condition is given as $\bm{u}(x,0) = \bm{u}^0(x)$. We know that
\[
	\bm{u} = \bm{K}\bm{w} \rightarrow \bm{w} = \bm{K}^{-1}\bm{u}
\]
The inversion of $\bm{K}$ is
\[
	\bm{K}^{-1} = \frac{1}{2a\rz} \begin{bmatrix}
		a & -\rz \\
		a & \rz
	\end{bmatrix},
\]
so we can compute
\[
	\bm{w}^0 = \bm{K}^{-1}\bm{u}^0 = \frac{1}{2a\rz} \begin{bmatrix}
		a & -\rz \\
		a & \rz
	\end{bmatrix}
	\begin{bmatrix}
		\rs^0 \\
		\us^0
	\end{bmatrix} =
	\frac{1}{2a\rz}
	\begin{bmatrix}
		a \rs^0 -\rz \us^0 \\
		a \rs^0 + \rz\us^0
	\end{bmatrix}
\]

\Q{Derive the solution in terms of decoupled characteristics}
\Q{Compute the solution in terms of original variables}

\end{document}
